\section{Week 13: Actual rANS Core Composition}

Week~13 closes \claimid{OBL-RANS-CORE-INVERSE} as
\status{proved, success-conditioned partial correctness}.  The result is a
direct composition theorem over three actual extracted procedures, not a
production full-HBZ wrapper refinement.

\subsection*{Actual execution path and theorem interface}

The existing \theoryname{Mode2RansActualInverse} harness executes:
\[
\begin{array}{c}
\procname{RansEncodeTarget.M._rans_encode}\\
\downarrow\\
\procname{HbzFullEncodeTarget.M.__copy_encoded_suffix}\\
\downarrow\\
\procname{RansDecodeTarget.M._rans_decode}.
\end{array}
\]
The copy and decoder calls occur only on the actual encoder's zero-
\texttt{bad} branch.  The compiled theorem
\path{actual_rans_encode_copy_decode_inverse} is a Hoare theorem over
\path{Mode2RansActualHarness.run}; it does not replace any of these calls
with a pure or observation procedure.

Its complete semantic precondition is exact initial argument binding and
\path{mode2_hbz_symbol_stream expected_symbols}.  In particular, it does
not assume encoder success, decoder reachability, copied trace equality,
decoder success, or decoded-symbol equality.

\subsection*{Encoder--copy interface}

On actual encoder success, the Week~11 theorem supplies
\[
\mathsf{segmentMatches}(encoded,off,trace)
\quad\land\quad off+|trace|=1024,
\]
where
\(trace=\mathsf{traceBytes}(\mathsf{symbolListOfArray}(symbols))\).
The actual copy theorem supplies
\(\mathsf{sliceEq}(encoded,copied,off,|trace|)\).
The Week~13 lemma \procname{copied_suffix_is_exact_trace} proves their
pointwise composition:
\[
\mathsf{segmentMatches}(copied,0,trace).
\]
Thus the decoder buffer is the actual post-state of the copy helper; it is not
an arbitrary witness.

\subsection*{Global trace to decoder byte reads}

The Week~12 decoder invariant phrases normalization reads by symbol index,
cursor, and intra-segment offset.  The new indexing lemma
\procname{trace_bytes_trace_segment_nth} proves that the local byte is the
corresponding element of the global trace:
\[
\mathsf{trace}[\mathsf{cursor}(i)+k]
 = \mathsf{segment}(i)[k].
\]
Combining this fact with the copied global segment yields
\path{segment_matches_implies_exact_decoder_segment_input}.  Its W64
corollary feeds every generated byte read required by the actual decoder
theorem.  No separate pointwise-read assumption remains in the harness
precondition.

\subsection*{W64/int size and configured state}

The actual harness computes
\[
encoded\_size = \mathsf{W64.of\_int}(1024)-encoder\_off.
\]
From the actual successful encoder bounds
\(0\leq \mathsf{uint}(off)\leq 1020\) and
\(\mathsf{uint}(off)+|trace|=1024\),
\procname{encoder_success_size_word_bridge} establishes
\[
\mathsf{uint}(encoded\_size)=|trace|,
\qquad
4\leq |trace|\leq 1024,
\]
and the exact W64 word equality.  The proof uses unsigned subtraction under a
proved order, rather than identifying modular and integer subtraction
implicitly.

The harness then performs three actual array stores.  The operation
\path{configured_decoder_state} denotes exactly that store sequence, and
\path{configured_decoder_state_fields} proves the resulting fields are
\((count,size,m)=(1024,|trace|,13)\).  Together with the copied segment, this
constructs the complete Week~12
\procname{actual_mode2_decoder_trace_input} predicate.

\subsection*{Failure/success disjunction}

The final theorem preserves both actual branches:
\[
\begin{array}{ll}
encoder\_bad\neq 0 &\Longrightarrow
  \neg decoder\_ran\ \land\ decoded=decoded_0\ \land\ state=state_0,\\[2mm]
encoder\_bad=0 &\Longrightarrow
  decoder\_ran\ \land\ decoder\_bad=0\ \land\\
&\quad decoder\_off=encoded\_size\ \land\\
&\quad decoded[0..1024)=symbols[0..1024)\ \land\\
&\quad decoded[1024..2048)=decoded_0[1024..2048).
\end{array}
\]
The success conclusion is obtained by sequentially applying the Week~11
encoder theorem, the actual copy theorem, and the Week~12 decoder theorem.
The generated loops are not unfolded again.

\subsection*{Partial correctness, non-vacuity, and rejected approaches}

The canonical-stream premise has a compiled all-zero array witness, so the
harness precondition is satisfiable.  The postcondition is nevertheless a
failure/success disjunction.  It does not establish that the success branch is
reachable, nor does it establish encoder or decoder termination/losslessness.
\claimid{OBL-RANS-ACTUAL-SUCCESS-WITNESS} therefore remains
\status{partial}.

A monolithic \texttt{inline; sim} proof was rejected because it duplicates
the completed encoder and decoder semantic invariants.  Supplying copied-trace
or decoder-read equality as a top-level premise was rejected as circular, and
constructing a decoder buffer witness was rejected because it severs the
actual copy-helper path.  The adopted decomposition uses only explicit
array-index, W64/int, and decoder-state adapters between the compiled actual
procedure theorems.

\subsection*{Security boundary and Week 14 gate}

This theorem is not the production
\procname{_encode_hb_z1_full}/\procname{_decode_hb_z1_full} inverse.  It proves
neither malformed-input rejection nor canonical parsing and does not justify
an encoding zero-loss or EUF-CMA claim.  Accordingly
\claimid{OBL-SIG-HBZ-ENCODE-DECODE} remains \status{partial}.

The Week~14 single gate is a success-conditioned composition at the two actual
full-HBZ wrappers, using the already proved prepare/apply inverse and the new
rANS core inverse.  Expansion to the \(h\) codec remains prohibited until that
wrapper theorem fresh-compiles.
